\section{需求分析}
\subsection{问题剖析}
灌溉决策的本质是在不确定条件下对“补水时机”和“补水量”作出选择。围绕这一问题，可以从以下四个维度分析其难点。

\textbf{需水的时间维度。}作物耗水由当天的辐射、温度、湿度与风速共同决定，具有明显的季节性和日际波动。在华北，5—7月参考蒸散可达每日6\,$\sim$\,9 mm，一次天气转折就可能使次日需水成倍变化。依据“昨天耗多少、今天补多少”的经验做法，在天气突变时会明显滞后。

\textbf{土壤的状态维度。}同样的次日需水，对根区水分充足的地块意味着无需灌溉，对接近胁迫阈值的地块则必须补水。因此，需水预测必须与根区水分亏缺一起使用，单看天气无法得出灌溉结论。FAO-56以“总可用水TAW—易利用水RAW—当前亏缺”描述这一状态，是业内通行的方法\cite{fao}。

\textbf{误差的风险维度。}预测偏低会推迟补水，造成作物水分胁迫；预测偏高会提前补水，增加无效用水和深层渗漏。两类错误的后果不对称，且与土壤缓冲能力、供水条件相关。只给出一个点预测，使用者无法判断风险大小，也就无法根据自身情况选择保守或节水的策略。

\textbf{约束的执行维度。}实际灌溉受机井出水量、管网压力、轮灌时段和劳动力等条件限制，计算出的“理论需水”未必能当天补足。若系统忽略供水能力，给出的建议可能无法执行，或让使用者误以为需水已经满足。

综合来看，核心痛点不在于缺少某一个预测模型，而在于“预测—误差—土壤—设备”之间缺少一条贯通、可检查的链条。这是本作品要解决的关键问题。

\subsection{需求调研}
本作品的需求调研以文献研究、公开数据分析和原型迭代三种方式开展。

\textbf{文献研究。}梳理FAO-56及其后续修订工作\cite{fao,pereira}，明确参考蒸散、作物系数、根区水量平衡和触发阈值的标准定义；调研亏缺灌溉相关研究\cite{fereres}与AquaCrop作物模型\cite{aquacrop}，明确水分胁迫与产量响应之间的关系及其所需的田间参数；调研国内ET\textsubscript{0}机器学习预测研究\cite{fan,feng}，确定树集成模型在有限气象输入下的适用性与常用评价指标。

\textbf{公开数据分析。}对三地2018—2025年逐日气象进行统计（图\ref{fig:climate}），确认参考蒸散在5—7月达到峰值、降水集中于7—8月，需水高峰与降水高峰存在错位。这说明春夏之交是灌溉决策压力最大、预测价值最高的时段，也决定了后续诊断要重点检查夏季月份的表现。

\textbf{原型迭代。}以典型使用流程为线索搭建原型，围绕“输入什么、看到什么、导出什么”多轮调整界面与接口。例如，早期版本只显示建议量，迭代后增加了“补水后预计亏缺”和“供水能力不足”提示；为防止页面状态与导出记录不一致，增加了服务端重算与请求编号校验。

在此基础上，确定目标用户为两类：一是需要比较不同灌溉情景、向农户提供建议的基层农技人员；二是农业水利、农业信息化相关课程的教师和学生。典型工作流程为：选择地区与日期，查看次日需水预测、误差区间和昨日值基线；输入当前根区亏缺、土壤总可用水、作物系数、灌溉效率和日供水上限；系统给出是否灌溉、毛灌溉量、净补水量与依据；使用者确认后导出本次决策记录，供后续复核。

\subsection{功能需求}
\begin{table}[H]\centering\small
\caption{功能需求与验收方式}\label{tab:requirements}
\begin{tabularx}{\linewidth}{@{}p{.18\linewidth}XX@{}}\toprule
需求 & 实现内容 & 验收方式\\\midrule
次日需水预测 & 基于目标日前49维特征推理次日ET\textsubscript{0} & 时间留出测试MAE优于昨日值基线\\
误差区间展示 & 曲线叠加名义90\%预测区间 & 独立年份校准，测试期检验覆盖率\\
基线对照 & 同屏显示昨日值与事后参考值 & 参考值仅用于核验，不进入模型输入\\
根区水量核算 & 显示剩余可用水、亏缺与触发线 & 参数变化时图形与计算同步更新\\
灌溉建议 & 输出毛灌溉量、净补水量与依据 & 与独立公式实现逐项一致\\
供水能力检查 & 建议量不超过日供水上限并提示 & 极小上限触发能力不足提示\\
记录导出 & JSON保存参数、策略、模型版本与结果 & 导出内容与页面当前输入一致\\
异常处理 & 非法参数、日期、地区返回错误 & 错误状态下禁止导出，修正后可恢复\\\bottomrule
\end{tabularx}
\end{table}

\subsection{非功能需求}
\begin{table}[H]\centering\small
\caption{非功能需求与实现方式}\label{tab:nonfunc}
\begin{tabularx}{\linewidth}{@{}p{.18\linewidth}X@{}}\toprule
维度 & 实现方式\\\midrule
可解释性 & 阈值、回填目标、效率与上限均以显式公式计算，界面同步展示依据\\
可复现性 & 保留原始接口响应与SHA-256哈希，固定依赖版本与随机种子\\
离线可用 & 模型、数据和页面部署在同一台计算机，推理不依赖外部服务\\
轻量部署 & 模型文件约416 KiB，单次推理中位耗时约11.5 ms，普通笔记本即可运行\\
数据安全 & 服务仅监听本机回环地址，不采集用户身份信息\\
跨端适配 & 手机竖屏布局适配，无横向溢出（局域网访问需另行配置服务地址）\\\bottomrule
\end{tabularx}
\end{table}

\subsection{评价指标设计}
预测精度、用水量和缺水风险是相互关联又彼此冲突的目标：更高的需水估计能降低缺水风险，却会增加用水；更低的估计则相反。为避免用任意加权的“综合得分”掩盖这一权衡，本作品分层评价：

\begin{itemize}
\item \textbf{模型层}：平均绝对误差MAE、均方根误差RMSE、决定系数$R^2$，以及预测区间的经验覆盖率与平均宽度；
\item \textbf{决策层}：季节毛灌溉量、缺水风险日（日末亏缺超过RAW的天数）、排水量和未满足需水量；
\item \textbf{系统层}：服务端参数校验、导出一致性、异常恢复与移动端适配。
\end{itemize}

这种分层设计使每一项结论都能对应到具体的数据与检验，也便于在接入田间数据后逐层替换和扩展。
